
\section{Predictors}
\begin{table}[ht]\centering\small
\caption{Predictors used by the boosted models. ``Mono.'' marks predictors on which PC-GBM is constrained to a non-decreasing response.}\label{tab:features}
\begin{tabular}{llll}\toprule
Predictor & Definition & Source & Mono.\\\midrule
slope & Slope angle (degrees) & GLO-30 & yes\\
slope\_max5 & Maximum slope in $5\times5$ window & GLO-30 & yes\\
relief510 & Max--min elevation in 510~m window & GLO-30 & yes\\
tpi510 & Elevation minus mean in 510~m window & GLO-30 & \\
relpos1k & Relative slope position in 1~km window & GLO-30 & \\
curv\_prof, curv\_plan & Profile and plan curvature (100/m) & GLO-30 & \\
slope\_sd3 & Std. dev. of slope in $3\times3$ window & GLO-30 & \\
log\_pga & $\log_{10}$ PGA (g) & ShakeMap & yes\\
log\_pgv & $\log_{10}$ PGV (cm\,s$^{-1}$) & ShakeMap & yes\\
log\_newmark & $\log_{10}(D_N+0.1)$, $D_N$ in cm & GLO-30 + ShakeMap & yes\\
\bottomrule\end{tabular}
\end{table}

\section{Covariate shift between events}
\begin{figure}[h]\centering
\includegraphics[width=\textwidth]{figS1_covariate_shift.pdf}
\caption{Distribution of (a) $\log_{10}$PGV, (b) slope and (c) $\log_{10}(D_N+0.1)$ over all domain cells (grey, weighted to the full domain) and over landslide cells (event colour). Thick bars span the interquartile range, thin bars the 5--95\% range, ticks mark medians.}\label{fig:shift}
\end{figure}

\section{Global logistic regression under low shaking}\label{sec:supp_lr}
Coefficients of the global LR (predictors $\log_{10}$PGV, slope, $\log_{10}$PGV$\times$slope, relief/100~m) fitted on the eight non-target events were $(-3.14, -0.066, 0.119, -0.74)$ with intercept $-2.17$ when Ays\'en was held out, $(-3.16, -0.173, 0.181, -0.65)$ with intercept $-2.20$ when Cinchona was held out, and $(-0.45, 0.045, 0.042, -0.75)$ with intercept $-6.01$ when Kaik\={o}ura was held out. The derivative of the log-odds with respect to slope, $\beta_{s}+\beta_{sv}\log_{10}\mathrm{PGV}$, becomes negative for PGV below $10^{0.066/0.119}\approx3.6$~cm\,s$^{-1}$ and $10^{0.173/0.181}\approx9.0$~cm\,s$^{-1}$ in the first two fits. The median ShakeMap PGV is 3.95~cm\,s$^{-1}$ over the Ays\'en domain and 10.1~cm\,s$^{-1}$ over the Cinchona domain, compared with 26--55~cm\,s$^{-1}$ in the other events. At Ays\'en, 95.4\% of cells have PGV below the first percentile of the training data.

\section{Gaussian tile random effect}\label{sec:supp_re}
In this alternative to the quasi-posterior of Eq.~(5) of the main text, the log-odds of cell $i$ in mapped tile $t$ are $m_i+b+u_t$ with $u_t\sim\mathcal{N}(0,\sigma_u^2)$. The tile likelihood is integrated over $u_t$ by an exact discrete Gaussian convolution of the tile log-likelihood on a grid of spacing 0.025 log-odds units; this was verified against brute-force numerical integration to eight significant digits. $\sigma_u$ and the conditional event offsets of the other events were estimated jointly by profile likelihood from their out-of-sample scores ($\hat\sigma_u=2.0$--$2.25$), and $\mu$ and $\tau$ of the prior were taken from these conditional offsets. The probability that an unmapped tile contains a landslide, needed for the selection term, was likewise averaged over $u_t$. Table~4 of the main text reports the results.

\section{Few-shot updating per event}
\begin{table}[ht]\centering\small
\caption{Per-event results of full HiPhyBoost updating (quasi-posterior with selection term): median $|\log_2 E/O|$ over 30 tile draws and coverage of the 90\% interval of the landslide total in the unmapped area. Niigata-Chuetsu has too few tiles for $k=40$.}\label{tab:perevent}
\begin{tabular}{lrrrrrrrr}\toprule
 & \multicolumn{4}{c}{$|\log_2 E/O|$} & \multicolumn{4}{c}{90\% coverage}\\\cmidrule(lr){2-5}\cmidrule(lr){6-9}
Event & $k$=1 & 5 & 10 & 40 & $k$=1 & 5 & 10 & 40\\\midrule
Kaik\={o}ura 2016 (NZ) & 0.65 & 0.47 & 0.51 & 0.36 & 0.90 & 0.77 & 0.57 & 0.63\\
Northridge 1994 (USA) & 0.50 & 0.55 & 0.46 & 0.43 & 1.00 & 0.87 & 0.80 & 0.73\\
Niigata-Chuetsu 2004 (JP) & 1.01 & 0.39 & 0.22 & -- & 0.60 & 0.63 & 0.67 & --\\
Iwate-Miyagi 2008 (JP) & 0.35 & 0.40 & 0.34 & 0.45 & 1.00 & 0.90 & 0.97 & 0.50\\
Ais\'{e}n 2007 (CL) & 0.68 & 0.45 & 0.34 & 1.15 & 1.00 & 1.00 & 1.00 & 0.80\\
Denali 2002 (USA) & 4.70 & 3.44 & 2.79 & 1.97 & 0.00 & 0.00 & 0.00 & 0.07\\
Cinchona 2009 (CR) & 3.01 & 0.78 & 0.48 & 0.64 & 0.20 & 0.60 & 0.73 & 0.23\\
Porgera 2018 (PNG) & 2.48 & 1.40 & 1.20 & 0.78 & 1.00 & 0.97 & 0.90 & 0.83\\
Palu 2018 (ID) & 0.63 & 0.53 & 0.66 & 0.52 & 0.90 & 0.93 & 0.87 & 0.53\\
\bottomrule\end{tabular}\end{table}


\section{Reliability diagrams}
\begin{figure}[h]\centering
\includegraphics[width=\textwidth]{fig4_reliability.pdf}
\caption{Reliability diagrams of transferred global LR and PC-GBM models for each held-out event (log--log axes; predictions grouped into 12 quantile bins). Points above the diagonal indicate underprediction, below the diagonal overprediction.}\label{fig:reliability}
\end{figure}

\section{Full metrics of Experiment 1}
{\small\begin{longtable}{llrrrrrr}
\caption{All performance metrics of Experiment~1. AP lift: average precision divided by prevalence; BSS: Brier skill score relative to a forecast that knows the true prevalence; $E/O$: expected over observed landslide cells; slope and intercept: logistic calibration of outcomes on predicted log-odds.}\label{tab:full}\\
\toprule Event & Model & AUC & AP lift & BSS & $E/O$ & Cal. slope & Cal. int.\\\midrule\endhead
Kaik\={o}ura 2016 (NZ) & Physics (Jibson 2000) & 0.883 & 9.59 & -1.637 & 8.28 & 0.36 & -3.61\\
Kaik\={o}ura 2016 (NZ) & Physics-logit & 0.861 & 9.20 & 0.007 & 1.5 & 2.19 & 5.81\\
Kaik\={o}ura 2016 (NZ) & Global LR & 0.826 & 9.91 & 0.009 & 1.63 & 1.53 & 2.12\\
Kaik\={o}ura 2016 (NZ) & GBM & 0.894 & 7.48 & 0.007 & 0.691 & 1.14 & 1.15\\
Kaik\={o}ura 2016 (NZ) & PC-GBM & 0.887 & 9.24 & 0.010 & 1.5 & 1.07 & -0.15\\
Kaik\={o}ura 2016 (NZ) & PC-GBM (local CV) & 0.934 & 16.79 & 0.017 & 1.03 & 0.90 & -0.46\\
\midrule
Northridge 1994 (USA) & Physics (Jibson 2000) & 0.715 & 3.61 & -0.361 & 1.88 & 0.22 & -3.56\\
Northridge 1994 (USA) & Physics-logit & 0.756 & 3.36 & 0.002 & 1.14 & 1.00 & -0.16\\
Northridge 1994 (USA) & Global LR & 0.723 & 3.94 & 0.002 & 1.44 & 1.03 & -0.17\\
Northridge 1994 (USA) & GBM & 0.788 & 2.43 & -0.029 & 1.04 & 0.43 & -3.00\\
Northridge 1994 (USA) & PC-GBM & 0.814 & 5.02 & 0.004 & 0.937 & 0.96 & -0.13\\
Northridge 1994 (USA) & PC-GBM (local CV) & 0.893 & 8.92 & 0.009 & 1.01 & 0.97 & -0.12\\
\midrule
Niigata-Chuetsu 2004 (JP) & Physics (Jibson 2000) & 0.763 & 2.96 & -0.127 & 1.18 & 0.22 & -2.75\\
Niigata-Chuetsu 2004 (JP) & Physics-logit & 0.767 & 2.94 & -0.005 & 0.175 & 1.50 & 4.74\\
Niigata-Chuetsu 2004 (JP) & Global LR & 0.750 & 3.27 & -0.002 & 0.197 & 0.82 & 0.64\\
Niigata-Chuetsu 2004 (JP) & GBM & 0.774 & 2.84 & -0.005 & 0.123 & 0.72 & 0.43\\
Niigata-Chuetsu 2004 (JP) & PC-GBM & 0.788 & 3.21 & -0.001 & 0.226 & 0.88 & 0.86\\
Niigata-Chuetsu 2004 (JP) & PC-GBM (local CV) & 0.809 & 3.19 & 0.010 & 1 & 0.73 & -1.03\\
\midrule
Iwate-Miyagi 2008 (JP) & Physics (Jibson 2000) & 0.777 & 4.06 & -1.295 & 7.76 & 0.40 & -3.59\\
Iwate-Miyagi 2008 (JP) & Physics-logit & 0.764 & 3.86 & 0.004 & 0.73 & 1.79 & 4.59\\
Iwate-Miyagi 2008 (JP) & Global LR & 0.750 & 4.56 & 0.008 & 0.978 & 0.96 & -0.16\\
Iwate-Miyagi 2008 (JP) & GBM & 0.796 & 4.58 & 0.008 & 1.19 & 1.00 & -0.17\\
Iwate-Miyagi 2008 (JP) & PC-GBM & 0.833 & 5.60 & -0.010 & 2.58 & 1.20 & -0.24\\
Iwate-Miyagi 2008 (JP) & PC-GBM (local CV) & 0.834 & 5.74 & 0.008 & 1.03 & 0.80 & -0.92\\
\midrule
Ais\'{e}n 2007 (CL) & Physics (Jibson 2000) & 0.609 & 1.66 & -27.005 & 82.9 & 0.11 & -6.69\\
Ais\'{e}n 2007 (CL) & Physics-logit & 0.739 & 2.53 & -0.318 & 16.8 & 1.55 & -0.37\\
Ais\'{e}n 2007 (CL) & Global LR & 0.247 & 0.60 & -0.068 & 10.8 & -0.81 & -13.05\\
Ais\'{e}n 2007 (CL) & GBM & 0.639 & 1.64 & -0.000 & 1.12 & 0.99 & -0.12\\
Ais\'{e}n 2007 (CL) & PC-GBM & 0.760 & 3.34 & 0.000 & 1.11 & 0.99 & -0.12\\
Ais\'{e}n 2007 (CL) & PC-GBM (local CV) & 0.826 & 3.90 & -0.000 & 1.04 & 0.98 & -0.11\\
\midrule
Denali 2002 (USA) & Physics (Jibson 2000) & 0.831 & 3.92 & -61.682 & 228 & 0.31 & -6.91\\
Denali 2002 (USA) & Physics-logit & 0.845 & 3.82 & -0.363 & 41.7 & 1.65 & -0.32\\
Denali 2002 (USA) & Global LR & 0.830 & 4.26 & -0.412 & 39.5 & 1.65 & -0.32\\
Denali 2002 (USA) & GBM & 0.850 & 4.65 & -0.328 & 21.8 & 1.55 & -0.29\\
Denali 2002 (USA) & PC-GBM & 0.805 & 3.44 & -0.201 & 23.4 & 1.53 & -0.26\\
Denali 2002 (USA) & PC-GBM (local CV) & 0.883 & 9.15 & 0.000 & 0.999 & 1.00 & -0.09\\
\midrule
Cinchona 2009 (CR) & Physics (Jibson 2000) & 0.559 & 1.47 & -0.246 & 0.934 & 0.13 & -3.33\\
Cinchona 2009 (CR) & Physics-logit & 0.734 & 2.52 & -0.001 & 0.31 & 0.78 & -0.12\\
Cinchona 2009 (CR) & Global LR & 0.397 & 0.74 & -0.004 & 0.365 & -0.40 & -7.51\\
Cinchona 2009 (CR) & GBM & 0.780 & 3.94 & -0.006 & 0.0184 & 0.86 & 2.84\\
Cinchona 2009 (CR) & PC-GBM & 0.774 & 2.94 & -0.006 & 0.0213 & 0.82 & 2.37\\
Cinchona 2009 (CR) & PC-GBM (local CV) & 0.812 & 4.96 & 0.008 & 0.982 & 0.81 & -0.79\\
\midrule
Porgera 2018 (PNG) & Physics (Jibson 2000) & 0.885 & 9.77 & -9.588 & 35.9 & 0.35 & -5.02\\
Porgera 2018 (PNG) & Physics-logit & 0.892 & 9.82 & -0.037 & 8.77 & 1.38 & -0.23\\
Porgera 2018 (PNG) & Global LR & 0.847 & 11.57 & -0.101 & 10.1 & 1.42 & -0.26\\
Porgera 2018 (PNG) & GBM & 0.876 & 8.55 & -0.052 & 5.56 & 1.29 & -0.21\\
Porgera 2018 (PNG) & PC-GBM & 0.872 & 9.05 & -0.097 & 8.84 & 1.40 & -0.25\\
Porgera 2018 (PNG) & PC-GBM (local CV) & 0.934 & 18.67 & 0.004 & 1.02 & 0.98 & -0.10\\
\midrule
Palu 2018 (ID) & Physics (Jibson 2000) & 0.742 & 3.00 & -1.974 & 8.69 & 0.19 & -4.58\\
Palu 2018 (ID) & Physics-logit & 0.762 & 3.03 & -0.001 & 2.07 & 1.10 & -0.19\\
Palu 2018 (ID) & Global LR & 0.733 & 3.34 & -0.006 & 2.44 & 1.13 & -0.20\\
Palu 2018 (ID) & GBM & 0.790 & 3.86 & 0.002 & 0.816 & 0.94 & -0.12\\
Palu 2018 (ID) & PC-GBM & 0.730 & 2.49 & -0.004 & 1.89 & 1.08 & -0.17\\
Palu 2018 (ID) & PC-GBM (local CV) & 0.838 & 4.61 & 0.002 & 0.97 & 0.96 & -0.12\\
\bottomrule
\end{longtable}
}

\section{The 38 earthquakes of the global extension}
{\small\begin{longtable}{lrrrcrrrrrl}
\caption{The 38 earthquakes of the global extension (ordered by year). Landslide cells and domain area on the 30~m grid; inventory type A = polygons (areas), P = points; ShakeMap stations; antecedent precipitation relative to its 1981--2010 climatology; leave-one-event-out AUC of PC-GBM and HiPhyBoost; HiPhyBoost event offset from the fit to all 38 events; source: ScienceBase item identifier of the inventory in the USGS repository (Schmitt et al. 2017; https://www.sciencebase.gov/catalog/item/<identifier>) or GNS Science (Kaik\={o}ura).}\label{tab:global}\\
\toprule Event & M & Landslide cells & km$^2$ & Inv. & Stations & Rain ratio & AUC PC-GBM & AUC HiPhy & $b_e$ & Source\\\midrule\endhead
San Fernando 1971 & 6.6 & 387 & 482 & P & 113 & 0.23 & 0.623 & 0.620 & -0.68 & \texttt{\scriptsize 5874cc26e4b0a829a320bbd4}\\
Guatemala 1976 & 7.5 & 6\,136 & 5\,265 & A & 0 & 0.59 & 0.691 & 0.698 & +0.23 & \texttt{\scriptsize 5874b3d4e4b0a829a320bb73}\\
Southern Italy 1980 & 6.9 & 68 & 11\,265 & P & 27 & 1.91 & 0.654 & 0.647 & -4.18 & \texttt{\scriptsize 5b6b2a35e4b006a11f779851}\\
Livermore 1980 & 5.8 & 105 & 908 & P & 12 & 1.20 & 0.683 & 0.697 & -0.54 & \texttt{\scriptsize 5874bdede4b0a829a320bbca}\\
Mammoth Lakes 1980 & 6.1 & 4\,024 & 2\,411 & A & 0 & 1.08 & 0.873 & 0.876 & +1.68 & \texttt{\scriptsize 5874bd42e4b0a829a320bbb9}\\
Coalinga 1983 & 6.7 & 3\,881 & 1\,856 & A & 48 & 2.37 & 0.791 & 0.798 & +1.76 & \texttt{\scriptsize 5874a164e4b0a829a320bb10}\\
San Salvador 1986 & 5.7 & 267 & 368 & P & 0 & 0.77 & 0.670 & 0.671 & +0.89 & \texttt{\scriptsize 5874ccaae4b0a829a320bbd9}\\
Loma Prieta 1989 & 6.9 & 1\,732 & 1\,905 & P & 101 & 2.41 & 0.686 & 0.694 & +0.21 & \texttt{\scriptsize 5874bd02e4b0a829a320bbb4}\\
Valle de la Estrella 1991 & 7.6 & 1\,643 & 2\,021 & A & 13 & 0.83 & 0.727 & 0.799 & -0.68 & \texttt{\scriptsize 5874bcc2e4b0a829a320bbaf}\\
Northridge 1994 (USA) & 6.7 & 10\,925 & 4\,377 & A & 185 & 0.41 & 0.809 & 0.831 & +1.38 & \texttt{\scriptsize 58595917e4b03639a6025f30}\\
Kobe 1995 & 6.9 & 2\,158 & 175 & A & 20 & 0.95 & 0.754 & 0.777 & +1.88 & \texttt{\scriptsize 58b538b7e4b01ccd54fddf59}\\
Central Italy 1997 & 6.0 & 176 & 9\,230 & P & 25 & 0.52 & 0.861 & 0.868 & -2.17 & \texttt{\scriptsize 5b6b075ee4b006a11f77930e}\\
Jueili 1998 & 5.7 & 847 & 1\,000 & A & 0 & 0.84 & 0.654 & 0.684 & +0.38 & \texttt{\scriptsize 5b5a20eee4b0610d7f4dcc52}\\
San Miguel 2001 & 7.7 & 134 & 8\,175 & P & 10 & 0.35 & 0.754 & 0.761 & -2.96 & \texttt{\scriptsize 5874cce7e4b0a829a320bbde}\\
Denali 2002 (USA) & 7.9 & 1\,579 & 15\,748 & A & 24 & 1.13 & 0.813 & 0.848 & -2.48 & \texttt{\scriptsize 5874a336e4b0a829a320bb15}\\
Lefkada 2003 & 6.3 & 252 & 223 & A & 9 & 1.10 & 0.864 & 0.878 & -0.02 & \texttt{\scriptsize 5874bc90e4b0a829a320bbaa}\\
Niigata-Chuetsu 2004 (JP) & 6.6 & 4\,860 & 410 & A & 448 & 1.17 & 0.799 & 0.800 & +2.45 & \texttt{\scriptsize 5874bd79e4b0a829a320bbbe}\\
Kashmir 2005 & 7.6 & 2\,929 & 4\,407 & A & 0 & 0.82 & 0.737 & 0.737 & -2.62 & \texttt{\scriptsize 5895dabfe4b0fa1e59bc788f}\\
Kiholo Bay 2006 & 6.7 & 379 & 229 & A & 23 & 0.97 & 0.872 & 0.864 & +0.14 & \texttt{\scriptsize 5874bb8de4b0a829a320bb9f}\\
Pisco 2007 & 8.0 & 271 & 21\,652 & P & 13 & 0.52 & 0.806 & 0.825 & -3.57 & \texttt{\scriptsize 5b5f3190e4b006a11f66e491}\\
Ais\'{e}n 2007 (CL) & 6.2 & 512 & 1\,247 & A & 0 & 0.74 & 0.654 & 0.691 & +0.84 & \texttt{\scriptsize 586fe5a7e4b01a71ba0c5f06}\\
Iwate-Miyagi 2008 (JP) & 6.9 & 4\,093 & 777 & A & 449 & 0.93 & 0.827 & 0.827 & +0.98 & \texttt{\scriptsize 5b59ec98e4b0610d7f4da850}\\
Wenchuan 2008 & 7.9 & 197\,160 & 77\,832 & A & 421 & 0.98 & 0.859 & 0.859 & +0.48 & \texttt{\scriptsize 586d7498e4b0f5ce109fc93b}\\
Cinchona 2009 (CR) & 6.1 & 4\,743 & 637 & A & 0 & 1.19 & 0.785 & 0.783 & +2.88 & \texttt{\scriptsize 61ead71bd34e8b818ad9f44b}\\
Haiti 2010 & 7.0 & 20\,702 & 3\,741 & A & 4 & 2.25 & 0.827 & 0.848 & +1.28 & \texttt{\scriptsize 585872dce4b0e40e53c241e1}\\
Tohoku-Oki 2011 & 9.1 & 3\,375 & 28\,934 & A & 1204 & 1.02 & 0.563 & 0.551 & -1.58 & \texttt{\scriptsize 5874a54be4b0a829a320bb2a}\\
Western Sichuan 2013 & 6.6 & 15\,278 & 5\,490 & P & 3 & 0.98 & 0.743 & 0.792 & +0.66 & \texttt{\scriptsize 59442c43e4b062508e3231be}\\
Gansu 2013 & 5.9 & 2\,169 & 320 & A & 2 & 1.28 & 0.705 & 0.655 & +2.64 & \texttt{\scriptsize 59442b04e4b062508e3231af}\\
Sichuan-Yunnan-Guizhou Region 2014 & 6.2 & 1\,024 & 363 & A & 0 & 1.20 & 0.747 & 0.772 & +0.91 & \texttt{\scriptsize 594428d4e4b062508e32319f}\\
Nepal 2015 & 7.8 & 20\,630 & 24\,517 & A & 2 & 1.08 & 0.795 & 0.822 & -0.80 & \texttt{\scriptsize 61f040e1d34e8b818adc3251}\\
Norcia 2016 & 6.6 & 349 & 4\,203 & P & 310 & 0.82 & 0.859 & 0.876 & -1.16 & \texttt{\scriptsize 5b6b2201e4b006a11f7795c3}\\
Kaik\={o}ura 2016 (NZ) & 7.8 & 28\,474 & 10\,832 & P & 227 & 1.05 & 0.892 & 0.917 & +0.62 & \texttt{\scriptsize GNS}\\
Visso 2016 & 6.1 & 204 & 673 & P & 87 & 0.82 & 0.755 & 0.800 & -0.46 & \texttt{\scriptsize 5b6b1177e4b006a11f7793ea}\\
Milin 2017 & 6.4 & 762 & 1\,755 & A & 0 & 0.94 & 0.770 & 0.754 & -1.14 & \texttt{\scriptsize 5f4d1e9c82ce4c3d12318d43}\\
Porgera 2018 (PNG) & 7.5 & 10\,455 & 24\,306 & A & 1 & 1.56 & 0.884 & 0.918 & -0.89 & \texttt{\scriptsize 61f04680d34e8b818adc32a8}\\
Palu 2018 (ID) & 7.5 & 7\,007 & 3\,665 & P & 0 & 1.04 & 0.729 & 0.776 & +0.47 & \texttt{\scriptsize 61f196a1d34e8b818add4fdf}\\
Lombok 2018 & 6.9 & 9\,265 & 1\,797 & A & 0 & 1.00 & 0.662 & 0.678 & +2.51 & \texttt{\scriptsize 5ef2514782ced62aaae3bcab}\\
Mesetas 2019 & 5.7 & 837 & 3\,101 & A & 0 & 1.47 & 0.950 & 0.927 & +0.66 & \texttt{\scriptsize 61ead1f5d34e8b818ad9f37f}\\
\bottomrule
\end{longtable}
}

\section{Zero-shot priors conditioned on event covariates}
\begin{table}[ht]\centering\small
\caption{Zero-shot priors for the event offset in the 38-event leave-one-event-out analysis. LPD: mean log predictive density of the held-out offset (nats); gain relative to the unconditional prior with paired Wilcoxon $p$ and permutation $p$ (200 random permutations of the covariate vectors among events); coverage of the 90\% interval for the landslide total; median interval width (ratio of bounds) and median error factor of the median forecast.}\label{tab:offreg}
\resizebox{\textwidth}{!}{\begin{tabular}{llrrrrrrr}\toprule
Covariates & Available & LPD & Gain & $p$ & $p_{\mathrm{perm}}$ & Coverage & Width & Error\\\midrule
None (unconditional prior) &  & -2.03 & -- & -- & -- & 33/38 & 364 & 2.9\\
F1: Mw, depth, mechanism, ShakeMap stations & forecast & -1.86 & +0.18 & 0.038 & 0.005 & 33/38 & 132 & 3.2\\
F2: PDSI, rain ratio, mean annual rain & forecast & -2.08 & -0.04 & 0.293 & 0.945 & 33/38 & 342 & 3.1\\
F3: F1 + F2 & forecast & -1.87 & +0.16 & 0.048 & 0.005 & 34/38 & 137 & 3.3\\
P1: inventory type, mapping extent, year & post-event & -1.99 & +0.04 & 0.153 & 0.030 & 33/38 & 245 & 3.6\\
P2: P1 + domain area & post-event & -1.83 & +0.20 & 0.010 & 0.005 & 36/38 & 120 & 3.4\\
P3: F3 + P2 & post-event & -1.85 & +0.18 & 0.067 & 0.005 & 35/38 & 121 & 3.3\\
\bottomrule\end{tabular}}\end{table}

